\documentclass{article}

% Use numeric citations (required for thebibliography without bibtex)
\PassOptionsToPackage{numbers,compress}{natbib}

% Workshop: single-blind reviewing
\usepackage[sglblindworkshop]{neurips_2026}
\workshoptitle{Agents in the Wild: Safety, Security, and Beyond (AIWILD @ NeurIPS 2026)}

\usepackage[utf8]{inputenc}
\usepackage[T1]{fontenc}
\usepackage{hyperref}
\usepackage{url}
\usepackage{booktabs}
\usepackage{amsfonts}
\usepackage{amsmath}
\usepackage{amssymb}
\usepackage{amsthm}
\usepackage{nicefrac}
\usepackage{microtype}
\usepackage{xcolor}
\usepackage{graphicx}
\usepackage{multirow}
\usepackage{listings}

\newtheorem{definition}{Definition}
\newtheorem{proposition}{Proposition}
\newtheorem{theorem}{Theorem}
\newtheorem{corollary}{Corollary}

\lstset{
  basicstyle=\ttfamily\small,
  breaklines=true,
  frame=single,
  xleftmargin=1em,
}

\title{Why LangGraph Cannot Hold a Fraud Block: Fail-Closed Orchestration for Agents with Consequential Authority}

\author{%
  Anonymous Author(s)
}

\begin{document}

\maketitle

\begin{abstract}
Agents deployed with real financial authority---booking travel, buying compute, paying
for APIs---are orchestrated by frameworks (LangGraph, AutoGen, OpenAI Agents SDK) that
cannot hold a safety decision. In a 15-line LangGraph program, a fraud gate that returns
\texttt{block} is silently overwritten by the next node's \texttt{allow} and the payment
executes. We show this is not a bug to be patched in application code but a structural
property of the framework's state semantics: any state channel whose reducer is
last-writer-wins is provably unable to be fail-closed (Corollary~\ref{cor:lww}), and we
characterize exactly what must replace it. Safe reducers are precisely the semilattice
operations that dominate the lattice join (Theorem~\ref{thm:general_char}); the join
itself---\texttt{block\_wins} on $\{\textit{allow}<\textit{flag}<\textit{block}\}$---is
the unique minimal one, so it is also the one with the fewest spurious escalations. A
delegation-closure result (Theorem~\ref{thm:delegation}) gives the single condition under
which sub-agents compose safely into a parent: the mapping between their escalation
vocabularies must be top-preserving.

We package these guarantees as \textbf{Commerce-Native Agentic Orchestration (CAO)}:
Commerce-Typed State channels, a payment-grounded state automaton with $O(n)$ validity
checking, Saga/2PC ceremony tiers, a structural atomic gate before the payment rail, a
confidence cascade, and a hash-chained tamper-evident audit log---together providing
Execution, Atomicity, Effectiveness, Verification, and Traceability (EAVT), none of which
existing frameworks provide at the semantic level. On ACP-Bench~\cite{acpbench} the
cascade reaches 88\% behavioral recall at 0\% hard false-positive rate across ten attack
classes at $<$\$0.01/1k sessions, versus 61\% for a full-session Claude Sonnet~4.6 judge at
\$3/1k. Code and results are open-source.
\end{abstract}

\section{Introduction}

LLM agents are transitioning from advisory roles to \emph{consequential actors}: systems authorized to book travel, subscribe to software, pay for API services, and execute wire transfers without per-action human approval. This transition demands that agent orchestration frameworks develop a formal account of what it means to take a financial action---not just what it means to route between graph nodes.

Distributed systems and payment rail engineering have provided formal answers to exactly these questions for four decades. The Saga pattern~\cite{saga} handles atomicity for long-running transactions. Two-phase commit~\cite{2pc} provides distributed consensus before irreversible commitment. ISO~8583~\cite{iso8583} defines the message lifecycle for financial transactions. Certificate Transparency~\cite{ct} provides tamper-evident audit logs via hash-chained, Merkle-anchored structures. These properties are not optional for commerce systems---they are load-bearing invariants. The problem is that no existing LLM orchestration framework enforces any of them.

\paragraph{Five transaction properties for agentic commerce.} We identify five properties that sound commerce orchestration must provide:

\begin{enumerate}
\item \textbf{Execution (E)}: the agent follows a valid state grammar $\mathcal{M}$; state-machine violations such as START$\to$COMMIT are caught before the payment rail executes.
\item \textbf{Atomicity (A)}: consequential actions commit all-or-nothing via risk-adjusted ceremony tiers (Saga at $L_1$, 2PC-lite at $L_2$); a RESERVE-then-failure produces a compensating VOID.
\item \textbf{Effectiveness (Ef)}: every consequential action passes through the Structural Atomic Gate before reaching the payment rail; no path exists from the reasoning layer to the rail that bypasses the gate.
\item \textbf{Verification (V)}: the four-tier confidence cascade (LocalRules $\to$ BehavioralML $\to$ LLM verifier $\to$ consensus) provides risk-adjusted verification depth.
\item \textbf{Traceability (T)}: the hash-chained CARLog produces a tamper-evident audit trail for every consequential action.
\end{enumerate}

\paragraph{The framework gap.} Table~\ref{tab:framework_compare} evaluates existing orchestration frameworks against these five properties. LangGraph, AutoGen, and OpenAI Agents were designed for graph routing and multi-agent coordination---none was designed to enforce the invariants that financial commitment requires.

\begin{table}[h]
\caption{Five transaction properties across orchestration frameworks. ``partial$^\dagger$'': provides a syntax-level or workflow mechanism but does not enforce the underlying safety semantic (see footnotes). CAO provides all five at the semantic level.}
\label{tab:framework_compare}
\centering
\footnotesize
\begin{tabular}{p{3.2cm}cccc}
\toprule
Property & LangGraph & AutoGen & OpenAI Agents & CAO \\
\midrule
Execution (grammar $\mathcal{M}$) & $\times$ & $\times$ & $\times$ & \checkmark \\
Atomicity (Saga/2PC) & $\times$ & $\times$ & $\times$ & \checkmark \\
Effectiveness (gate) & partial$^a$ & $\times$ & partial$^b$ & \checkmark \\
Verification (cascade) & $\times$ & $\times$ & partial$^b$ & \checkmark \\
Traceability (CARLog) & partial$^c$ & $\times$ & $\times$ & \checkmark \\
\bottomrule
\multicolumn{5}{p{11cm}}{\footnotesize
$^a$ \texttt{interrupt\_before} gates workflow execution but provides no fraud-semantic protocol-validity check; LWW semantics allow a subsequent node to overwrite a block (Theorem~\ref{thm:lww_unsound}).
$^b$ OpenAI Agents SDK guardrails observe input/output text; they enforce no grammar over the financial action sequence and produce no tamper-evident audit log.
$^c$ LangSmith provides observability logging; it does not provide hash-chain tamper-evidence or Merkle inclusion proofs verifiable by third parties.}
\end{tabular}
\end{table}

\paragraph{The three-layer stack.} Agentic commerce safety operates at three distinct layers:

\begin{itemize}
\item \textbf{Layer~0 (Agent Reasoning)}: the full agentic stack---LLM, tools, memory, planning loop---generates consequential action calls. Adversaries inject into prompts or manipulate reasoning. Safety tools at this layer (garak, promptfoo, LlamaGuard, FinHarness) observe agent text output, conversation content, and tool selection decisions. They cannot observe the protocol wire or wallet state below.

\item \textbf{Layer~1 (MCP/Protocol, this paper)}: the agent's structured \texttt{ConsequentialAction} messages traverse the MCP wire before reaching any payment rail. The wire carries typed fields (action type, merchant, MCC, amount, idempotency key, session context, wallet state) that Layer~0 tools cannot observe. Adversaries exploit the payment-protocol state machine---issuing \texttt{COMMIT} before \texttt{QUOTE}, replaying idempotency keys, submitting out-of-allowlist MCCs, exceeding spend ceilings. CAO operates here, enforcing all five EAVT properties at the protocol wire.

\item \textbf{Layer~2+ (Payment Rail)}: Stripe Radar, payment gateway velocity checks, issuing bank authorization, merchant fraud systems. These operate \emph{after} the agent has called \texttt{payment\_intents.create()} or equivalent---the agent is already deep into the transaction. Layer~2+ can decline a charge, but it cannot provide graceful in-session recovery, issue a compensating void to an earlier reservation, or access the LLM session history that led to the call.
\end{itemize}

The three layers are not competing---they are defense-in-depth. Layer~0 is the correct place to catch a jailbroken agent. Layer~1 is the correct place to catch a correctly-reasoning agent whose \emph{action} violates payment-protocol invariants. Layer~2+ is the backstop for attacks that evade both. This paper formalizes Layer~1.

\paragraph{Implications for deployed agent frameworks.}
Practitioners today defend fraud gates by convention: ``don't write to
\texttt{decision} after the gate.'' Theorem~\ref{thm:general_char} shows why convention
is insufficient and what the fix is. Because safety is a property of the reducer algebra,
it can be enforced once, at channel-declaration time, and no downstream node---including
an LLM-driven one whose outputs are not under the developer's control---can regress it.
The same theorem tells framework maintainers that offering \texttt{block\_wins} (or any
join) as a first-class channel type is the \emph{minimal} change that makes their framework
able to hold a safety decision, and that offering anything stricter only increases
false escalations. Theorem~\ref{thm:delegation} gives the one rule for multi-agent
deployments: when a sub-agent reports back, map its ``blocked'' to the parent's
``blocked''---anything else silently drops the sub-agent's veto. Proposition~\ref{prop:decidability}
adds that protocol validity of the resulting action stream is checkable in linear time
without any LLM call, so the whole gate runs inline at $<$10\,ms.

\paragraph{What is novel here.}
None of the individual CAO components is new: DFA simulation, Saga/2PC, CRDT joins,
and hash-chain logs each have decades of literature. The contribution is:
(1)~a framework-independent characterization (Theorem~\ref{thm:general_char}) of
\emph{exactly} which state-channel reducers can carry a safety decision---safe reducers
are the semilattice operations dominating the lattice join, and the join is the unique
minimal one---so that fail-closed safety is a property of the reducer algebra, not of
application code; (2)~a delegation-closure theorem (Theorem~\ref{thm:delegation}) giving
the necessary and sufficient condition under which hierarchical agent populations remain
fail-closed; (3)~identifying EAVT as the property set that consequential-action
orchestration must provide and proving that the CAO composition provides all five while
existing frameworks provide none at the semantic level. These results formally explain
why LangGraph's last-writer-wins semantics are a structural safety gap rather than a
configuration choice, and why \texttt{block\_wins} is not one design among many but the
precision-optimal safe reducer.

\paragraph{Contributions.} Each contribution maps to one or more of the five transaction properties:

\begin{enumerate}
\item \textbf{Commerce State Automaton} $\mathcal{M}$ (\S\ref{sec:automaton}): a formal finite automaton over ISO~8583-grounded action states whose acceptance language defines the security invariant for agentic payment sessions. Provides the \textbf{Execution} property; decidability and $O(n)$ complexity proven.

\item \textbf{Safe reducer theory and Commerce-Typed State (CTS)} (\S\ref{sec:general_typing}--\ref{sec:cts}): a lattice-theoretic characterization of safe state-channel reducers (Theorem~\ref{thm:general_char}), its corollaries (last-writer-wins is unsafe on every lattice; the join is precision-optimal), and delegation closure for agent hierarchies (Theorem~\ref{thm:delegation}). CTS instantiates this on $\{\textit{allow}<\textit{flag}<\textit{block}\}$ with \texttt{block\_wins}; Theorems~\ref{thm:type_safety}--\ref{thm:char} give the concrete type-safety, LWW-unsoundness, and minimality statements. Provides the \textbf{Verification} and \textbf{Effectiveness} properties at the state level.

\item \textbf{Structural Atomic Gate} (\S\ref{sec:gate}): formalizes what LangGraph's \texttt{interrupt\_before} cannot provide; the Commerce Action Receipt (CAR) abstraction achieves 2PC-style atomicity. Provides the \textbf{Effectiveness} property.

\item \textbf{Ceremony Tiers} (\S\ref{sec:tiers}): a principled mapping from transaction risk to commitment protocol (Saga at $L_1$, 2PC-lite at $L_2$, consensus at $L_3$). Provides the \textbf{Atomicity} property.

\item \textbf{CARLog} (\S\ref{sec:carlog}): hash-chained, Merkle-anchored audit log producing tamper-evident provenance for every consequential action. Provides the \textbf{Traceability} property.

\item \textbf{Property evaluation} (\S\ref{sec:experiments}): measured validation of Execution, Effectiveness, and Traceability across 900 sessions; Atomicity and Verification are architectural guarantees proven by construction.
\end{enumerate}

\section{Background and Related Work}

\subsection{LLM Agent Orchestration Frameworks}

\textbf{ReAct}~\cite{react} alternates reasoning trace generation with tool execution. No state management beyond context window; no safety semantic on tool calls.

\textbf{LangGraph}~\cite{langgraph}: stateful graph execution based on Pregel~\cite{pregel}. State is a typed dictionary with per-channel reducers; \texttt{interrupt\_before} pauses for human review. Key limitation: reducers are user-specified---the framework imposes no constraints on reducer semantics. \texttt{interrupt\_before} is a workflow gate, not a fraud gate.

\textbf{AutoGen}~\cite{autogen}: multi-agent conversational actors. Safety mechanism: role-based prompt engineering. Does not model financial action lifecycle.

\textbf{OpenAI Agents SDK}~\cite{openai_agents}: guardrails on agent inputs and outputs. Guardrails observe text, not payment-rail state. Partial effectiveness and verification arise from guardrail checks, but these checks enforce no grammar over the financial action sequence and produce no audit log.

\textbf{Summary}: all existing frameworks operate at the reasoning and routing layer. None provides any of the five EAVT properties for financial commitments.

\subsection{Payment Protocol Foundations}

\textbf{ISO 8583}~\cite{iso8583}: core message types enforce ordering---0100 (authorization request) $\to$ 0200 (financial transaction) $\to$ 0220 (financial advice) $\to$ 0420 (reversal). This is precisely the invariant that cold-start attacks (B1) violate.

\textbf{x402 AI-payment protocol}~\cite{x402}: maps the ISO~8583 lifecycle to HTTP headers for AI-to-AI payments, and documents the revert-grant attack class (B9) that occurs when agents can issue \texttt{REFUND} after \texttt{SETTLE}.

\textbf{Idempotency in payment rails}: Stripe, Adyen, and ISO~20022 require idempotency keys on all mutation requests. If the agent verifier lacks access to the settled-key registry, replay attacks (B5) are invisible.

\textbf{The agentic gap}: in all prior payment systems, the ISO~8583 state machine was enforced by terminal hardware and issuing bank logic. An LLM agent issuing structured API calls directly to a payment rail can skip states. \textbf{There is no terminal. The agent is the terminal.}

\subsection{Distributed Transaction Theory}

\textbf{Two-Phase Commit}~\cite{2pc}: coordinator asks all participants to prepare (phase~1); commits if all agree, aborts otherwise (phase~2). Guarantees no participant commits if any participant aborts.

\textbf{Saga Pattern}~\cite{saga}: sequence $T_1, \ldots, T_n$ where each $T_i$ has compensating transaction $C_i$. Failure at $T_j$ triggers $C_{j-1}, \ldots, C_1$. Trades atomicity for availability.

\textbf{CRDTs}~\cite{crdts}: join-semilattice data structures whose operations satisfy commutativity, associativity, and idempotency. The join operation \emph{encodes the invariant}---a G-Counter's join is \texttt{max} (can only increase). The algebraic structure of the join \emph{is} the safety guarantee. Our Theorem~\ref{thm:general_char} is the converse direction for safety channels: it shows that the CRDT algebra is not merely sufficient but \emph{forced} by natural operational safety requirements.

\textbf{Behavioral and session types.} Typestate~\cite{typestate} and session types~\cite{sessiontypes} constrain \emph{which operations} a value or channel may undergo next. CTS is a complementary discipline on \emph{which values} a channel may take next: the reducer type restricts the channel's evolution to a monotone path in a safety lattice. Where session types reject ill-ordered protocol messages at compile time, CTS rejects safety-regressing writes; both encode a behavioral invariant in the type rather than in application logic.

\section{Commerce-Native Agentic Orchestration (CAO)}
\label{sec:cao}

\subsection{General Safety Typing for State Channels}
\label{sec:general_typing}

We first state the framework-independent result. A graph-execution framework
(Pregel~\cite{pregel}, LangGraph~\cite{langgraph}) exposes \emph{state channels}:
a value domain together with a reducer that folds node writes into the channel.
We ask which reducers can carry a safety decision at all.

\begin{definition}[Bounded safety lattice]
A \emph{bounded safety lattice} is a finite lattice $(L,\leq)$ with least element
$\bot$ (``no concern'') and greatest element $\top$ (``blocked''). We write
$\sqcup$ for its join.
\end{definition}

\begin{definition}[Channel reducer]
A \emph{channel reducer} on $L$ is a map $r: L\times L\to L$ with $r(\bot,x)=x$
(an empty channel adopts its first write). Given node writes $b_1,\dots,b_n\in L$,
the channel state is the fold $D_0=\bot$, $D_i=r(D_{i-1},b_i)$.
\end{definition}

\begin{definition}[Safe reducer]
\label{def:safe}
$r$ is \emph{safe} for $L$ if for every finite trace $b_1,\dots,b_n$:
\begin{enumerate}
\item[(S1)] \emph{Fail-closed}: $b_i=\top \Rightarrow D_j=\top$ for all $j\ge i$.
\item[(S2)] \emph{Non-regressive}: $D_{i-1}\le D_i$ and $b_i\le D_i$ for all $i$.
\item[(S3)] \emph{Order-independent}: every permutation of the trace yields the same $D_n$
  (concurrent writes within a superstep fold deterministically).
\item[(S4)] \emph{Replay-stable}: duplicating any $b_i$ leaves $D_n$ unchanged
  (idempotent retries).
\end{enumerate}
\end{definition}

\begin{theorem}[Characterization of safe reducers]
\label{thm:general_char}
Let $L$ be a bounded safety lattice. A channel reducer $r$ is safe for $L$ if and
only if $r$ is a semilattice operation (associative, commutative, idempotent)
that dominates the join: $x\sqcup y\le r(x,y)$ for all $x,y\in L$.
Consequently $\sqcup$ is the unique \emph{pointwise-minimal} safe reducer:
every other safe reducer satisfies $r(x,y)>x\sqcup y$ for some pair, i.e.\ it
escalates strictly more often than $\sqcup$ without adding safety.
\end{theorem}

\begin{proof}[Proof sketch; full proof in Appendix~\ref{app:general_char}]
($\Leftarrow$) A semilattice operation satisfies (S3) and (S4); $x\sqcup y\le r(x,y)$
gives (S2); (S1) follows from (S2) because $\top$ is maximal.
($\Rightarrow$) Using $r(\bot,x)=x$: (S3) on two- and three-element traces yields
commutativity and associativity; (S4) on the trace $(x,x)$ yields idempotence;
(S2) yields $x\le r(x,y)$ and $y\le r(x,y)$, hence $x\sqcup y\le r(x,y)$ since
$x\sqcup y$ is the least upper bound. Minimality: $\sqcup$ is itself safe and is
pointwise $\le$ every safe $r$. \qed
\end{proof}

\begin{corollary}[Chains; \texttt{block\_wins}]
\label{cor:chain}
If $L$ is a chain then $\sqcup=\max$ is the unique minimal safe reducer. For
$L=\{\textit{allow}<\textit{flag}<\textit{block}\}$, $\max$ is exactly \texttt{block\_wins}.
\end{corollary}

\begin{corollary}[Last-writer-wins is unsafe on every non-trivial lattice]
\label{cor:lww}
$r(x,y)=y$ violates (S2) whenever $|L|\ge2$: $r(\top,\bot)=\bot\not\ge\top$.
No framework whose default channel semantics is last-writer-wins can carry a
fail-closed decision without replacing the reducer.
\end{corollary}

\begin{corollary}[Precision optimality]
\label{cor:precision}
Among safe reducers, $\sqcup$ has an empty spurious-escalation set
$\{(x,y): r(x,y)>x\sqcup y\}$. Any stricter safe reducer only adds escalations
(soft-flag false positives), never safety. This is the formal reason the empirical
escalation rate in \S\ref{sec:experiments} is attributable to the verifiers, not the reducer.
\end{corollary}

\begin{theorem}[Delegation closure]
\label{thm:delegation}
Let a parent graph $P$ carry a safe reducer $r_P$ on $L_P$ and delegate to a
sub-agent $S$ carrying a safe reducer $r_S$ on $L_S$; $S$'s final state $d^S$ is
written back into $P$'s channel through an embedding $\iota:L_S\to L_P$.
The composite is fail-closed---(i) a block in $P$ before delegation persists
regardless of $d^S$, and (ii) a block in $S$ forces $P$ to $\top_P$---if and only
if $\iota(\top_S)=\top_P$. If moreover $\iota$ is monotone, severity order is
preserved across the boundary.
\end{theorem}

\begin{proof}
(i) $D^P_{i-1}=\top_P\Rightarrow D^P_i=r_P(\top_P,\iota(d^S))=\top_P$ by (S1).
(ii) If $d^S=\top_S$ then $\iota(d^S)=\top_P$ and $D^P_i=\top_P$ by (S1).
Conversely, if $\iota(\top_S)=t<\top_P$, the trace with $D^P_{i-1}=\bot$ yields
$D^P_i=t\ne\top_P$: the sub-agent's block is lost. Monotone $\iota$ gives
$a\le_S b\Rightarrow\iota(a)\le_P\iota(b)$, so (S2) in $S$ is reflected in $P$. \qed
\end{proof}

Theorem~\ref{thm:delegation} is the compositional-safety statement: an agent
population in which every member uses a safe reducer and every delegation edge is
top-preserving is fail-closed as a whole, regardless of depth or fan-out. The
remainder of this section instantiates the framework for agentic commerce with
$L=\{\textit{allow}<\textit{flag}<\textit{block}\}$.

\subsection{Execution Property: Commerce State Automaton}
\label{sec:automaton}

\begin{definition}[Commerce State Automaton]
Let $\mathcal{M} = (Q, \Sigma, \delta, q_0, F)$ where:
\begin{itemize}
\item $Q = \{\texttt{START}, \texttt{FIND}, \texttt{QUOTE}, \texttt{RESERVE}, \texttt{COMMIT}, \texttt{SETTLE}, \texttt{VOID}, \texttt{END}\}$
\item $\Sigma = \{\textit{find}, \textit{quote}, \textit{reserve}, \textit{commit}, \textit{settle}, \textit{void}, \textit{refund}\}$
\item $q_0 = \texttt{START}$, $F = \{\texttt{END}\}$
\item $\delta$ defines valid transitions per Table~\ref{tab:transitions}
\end{itemize}
\end{definition}

\begin{table}[h]
\caption{Valid transitions in $\mathcal{M}$ with ISO~8583, x402, and Stripe PaymentIntent grounding.}
\label{tab:transitions}
\centering
\small
\begin{tabular}{llll}
\toprule
From & Action & To & Protocol Grounding \\
\midrule
\texttt{START} & find & \texttt{FIND} & Session initialization \\
\texttt{FIND} & find & \texttt{FIND} & Iterative discovery (ISO~8583 pre-auth inquiry loop) \\
\texttt{FIND} & quote & \texttt{QUOTE} & Proceed to pricing \\
\texttt{QUOTE} & reserve & \texttt{RESERVE} & Hold funds---2PC prepare phase \\
\texttt{QUOTE} & commit & \texttt{COMMIT} & Direct commit (Saga, low-risk) \\
\texttt{RESERVE} & commit & \texttt{COMMIT} & Execute after hold---2PC commit phase \\
\texttt{RESERVE} & void & \texttt{VOID} & Cancel hold---Saga compensating transaction \\
\texttt{COMMIT} & settle & \texttt{SETTLE} & ISO~8583 0220 financial advice \\
\texttt{COMMIT} & refund & \texttt{VOID} & Saga compensating transaction \\
Any & --- & \texttt{END} & Session close \\
\bottomrule
\end{tabular}
\end{table}

\begin{definition}[Valid Session]
A session trajectory $s = a_1, a_2, \ldots, a_n$ is \emph{protocol-valid} iff:
\begin{enumerate}
\item The state sequence $q_0, \delta(q_0, a_1), \delta(q_1, a_2), \ldots \in F$ is accepted by $\mathcal{M}$ (state-transition validity)
\item For all $a_i$ with $a_i.\textit{action} = \textit{commit}$: $a_i.\textit{amt} \leq \textit{wallet.per\_transaction}$ (field validity---spend limit)
\item For all $a_i$ with $a_i.\textit{action} = \textit{commit}$: $a_i.\textit{mcc} \in \textit{wallet.allowed\_mcc}$ (field validity---MCC allowlist)
\item All idempotency keys $k$ in committed actions are unique in the settled-key registry (field validity---idempotency)
\item $\textit{ctx.agent\_id}$ matches the registered agent for $\textit{ctx.session\_id}$ (identity validity)
\end{enumerate}
\end{definition}

\textbf{Every payment-protocol attack in the taxonomy is a violation of exactly one of these five conditions.} B1 violates condition~1; B5 violates condition~4; B7 violates condition~3; B8 violates condition~2; B4 violates condition~5.

\begin{proposition}[Decidability and complexity of protocol validity]
\label{prop:decidability}
Let $s = a_1, \ldots, a_n$ be a session trajectory of length $n$. Then:
\begin{enumerate}
\item[(a)] $L(\mathcal{M})$ (the set of state-transition valid trajectories) is a regular language over $\Sigma^*$. Its complement $\overline{L(\mathcal{M})}$ is also regular (regular languages are closed under complement).
\item[(b)] Membership testing---``is $s$ state-transition valid?''---is decidable in $O(n)$ time and $O(|Q|)$ space by single-pass DFA simulation.
\item[(c)] Conditions 2--5 (field and identity validity) are decidable in $O(n)$ time given access to wallet state, the settled-key registry, and the agent-ID registry.
\item[(d)] Full protocol validity (all five conditions) is therefore decidable in $O(n)$ time. The protocol-violation decision problem is in $\mathbf{P}$.
\end{enumerate}
\end{proposition}

\begin{proof}
(a): $\mathcal{M}$ is a finite automaton; its acceptance language is regular by definition (Kleene, 1956). Closure under complement holds for all regular languages (DFA complement via state complementation). (b): DFA simulation visits each input symbol once and maintains a single state pointer over $|Q| = 8$ states---$O(n)$ time, $O(1)$ state, $O(|Q|)$ total. (c): Conditions 2 and 3 require one comparison per commit action against wallet thresholds (constant-time lookups); condition 4 requires one set-membership query per idempotency key (hash-set, $O(1)$ amortized); condition 5 requires one registry lookup per session. All $O(n)$ total. (d): Sequential composition of $O(n)$ procedures is $O(n)$; $O(n) \subset \mathbf{P}$. \qed
\end{proof}

\begin{corollary}[Protocol-violation certificate]
A protocol-invalid trajectory $s$ has a polynomial-time checkable certificate: the first action $a_i$ violating any of the five conditions. A verifier can produce and the payment rail can verify this certificate in $O(n)$ time, enabling lightweight dispute resolution for chargebacks.
\end{corollary}

\subsection{Verification and Effectiveness: Commerce-Typed State}
\label{sec:cts}

LangGraph state channels are (type, reducer) pairs. The default reducer is last-writer-wins. We observe that the reducer can encode a safety invariant as a monotone semilattice join.

\begin{definition}[Commerce-Typed Channel]
A \emph{commerce-typed channel} is a (type, reducer) pair where the reducer $f$ is a monotone semilattice join over the value domain:
\begin{itemize}
\item Associative: $f(f(x,y),z) = f(x,f(y,z))$
\item Commutative: $f(x,y) = f(y,x)$
\item Idempotent: $f(x,x) = x$
\item Monotone: $x \leq f(x,y)$ for all $y$
\end{itemize}
\end{definition}

\begin{proposition}[CTS channels are CRDTs]
Every CTS channel is a join-semilattice CRDT. Any two replicas receiving the same updates converge to the same state regardless of order; the channel state is a lower bound that can only grow more committed, never less.
\end{proposition}

\textbf{Concrete CTS channels:}
\begin{lstlisting}
spend:      Annotated[Decimal, operator.add]      # monotone: spend_new >= spend_old
decision:   Annotated[str, block_wins]             # block > flag > allow
provenance: Annotated[list[CARRef], operator.add]  # append-only audit trail
flags:      Annotated[frozenset, frozenset.union]  # flags never removed
\end{lstlisting}

where \texttt{block\_wins(a, b) = "block" if "block" in \{a,b\} else "flag" if "flag" in \{a,b\} else "allow"}.

\begin{proposition}[Safety monotonicity]
\label{thm:monotone}
In a CTS-typed graph state, if any verifier produces decision=\texttt{"block"} for any action in a session, the session-level decision is irreversibly \texttt{"block"} regardless of subsequent verifier outputs.
\end{proposition}

\begin{proof}[Proof sketch]
By definition of \texttt{block\_wins}, $f(x, \texttt{"block"}) = \texttt{"block"}$ for all $x \in \{\texttt{allow}, \texttt{flag}, \texttt{block}\}$. Idempotency: $f(\texttt{"block"}, \texttt{"block"}) = \texttt{"block"}$. Monotonicity: once the lattice element reaches \texttt{"block"}, no join can reduce it. \qed
\end{proof}

\begin{corollary}[Fail-closed safety]
The \texttt{block\_wins} reducer provides formal \emph{fail-closed safety}---the agent pipeline cannot unblock itself by producing subsequent allow decisions. LangGraph's last-writer-wins scalar semantics allows exactly this: a subsequent node can overwrite a prior block if the state is not carefully protected.
\end{corollary}

\begin{proposition}[Rules-certainty veto under CTS]
\label{prop:veto}
The protocol-layer veto (block$\to$flag downgrade when LocalRules score$=0$ and ML uncertainty$>0.20$) does not violate Theorem~\ref{thm:monotone}. The veto fires at the \emph{action level}---it modifies the ACP node's output $b_i$ from \texttt{block} to \texttt{flag} before the session-level reducer runs. Crucially, the veto condition requires LocalRules decision$=\texttt{"allow"}$ (score$=0$), which means LocalRules has not previously issued a block for this action. The session-level decision after step $i$ is $D_i = \texttt{block\_wins}(D_{i-1}, b_i)$. If $D_{i-1} = \texttt{block}$ (a prior action was blocked), then $D_i = \texttt{block\_wins}(\texttt{block}, \texttt{flag}) = \texttt{block}$---the veto cannot reduce a decision already in the log. The veto is a precision mechanism on uncertain ML outputs, not a safety bypass.
\end{proposition}

\begin{theorem}[CTS Type Safety]
\label{thm:type_safety}
Let $G$ be a CTS-typed graph where all decision channels carry reducer \texttt{block\_wins}. Let $T = (s_0, s_1, \ldots, s_n)$ be any execution trace of $G$. If there exists $i \in \{0,\ldots,n\}$ such that the ACP node at step $i$ outputs $b_i = \texttt{block}$, then for all $j \geq i$: $\mathrm{decision}(s_j) = \texttt{block}$.
\end{theorem}

\begin{proof}
By induction on $j - i$.

\textit{Base case} ($j = i$): $\mathrm{decision}(s_i) = \texttt{block\_wins}(\mathrm{decision}(s_{i-1}), b_i) = \texttt{block\_wins}(d, \texttt{block})$ for some $d \in D$. By the definition of \texttt{block\_wins}, for all $d \in \{\texttt{allow}, \texttt{flag}, \texttt{block}\}$: $\texttt{block\_wins}(d, \texttt{block}) = \texttt{block}$. So $\mathrm{decision}(s_i) = \texttt{block}$.

\textit{Inductive step} ($j = k{+}1$, assuming $\mathrm{decision}(s_k) = \texttt{block}$): $\mathrm{decision}(s_{k+1}) = \texttt{block\_wins}(\mathrm{decision}(s_k), b_{k+1}) = \texttt{block\_wins}(\texttt{block}, b_{k+1})$. By \texttt{block\_wins} definition, this equals \texttt{block} for all $b_{k+1} \in D$. So $\mathrm{decision}(s_{k+1}) = \texttt{block}$.

By induction, $\mathrm{decision}(s_j) = \texttt{block}$ for all $j \geq i$. \qed
\end{proof}

\begin{theorem}[LWW Unsoundness]
\label{thm:lww_unsound}
Let $G$ be an execution graph where the decision channel uses last-writer-wins: $r(x, y) = y$. Then for every implementation of a fraud gate node $n_1$ that can output \textup{\texttt{block}}, there exists a two-node execution trace $(n_1 \to n_2 \to \mathrm{payment\_rail})$ in which $n_1$ outputs \textup{\texttt{block}} and the payment rail executes.
\end{theorem}

\begin{proof}
Construct $n_2$ such that it outputs $b_2 = \texttt{allow}$ (e.g., a reasoning node that re-evaluates the action positively, or an identity node that passes through a cached allow decision). With LWW: $\mathrm{decision}(s_2) = r(\texttt{block}, \texttt{allow}) = \texttt{allow}$. The payment rail routing condition is $\mathrm{decision} \neq \texttt{block}$, which is satisfied; the payment rail executes. The node $n_2$ is a standard construction in LangGraph (any graph node whose output includes \texttt{decision: "allow"} overwrites the channel value). Therefore, for any fraud gate $n_1$, the two-node trace $n_1 \to n_2$ bypasses the block. \qed
\end{proof}

\begin{corollary}
The counterexample in \S\ref{sec:gate} is not a pathological edge case but an instance of Theorem~\ref{thm:lww_unsound}: any LangGraph graph using LWW on a \texttt{decision: str} channel admits this construction. Adding CTS with \texttt{block\_wins} is necessary (Theorem~\ref{thm:char} below), not merely sufficient.
\end{corollary}

\begin{theorem}[Commerce instantiation of Theorem~\ref{thm:general_char}]
\label{thm:char}
On $D=\{\texttt{allow}<\texttt{flag}<\texttt{block}\}$, a reducer is safe in the sense of
Definition~\ref{def:safe} iff it is a semilattice operation dominating $\max$; \texttt{block\_wins}
$=\max$ is the unique minimal such reducer. Hence CTS with \texttt{block\_wins} is the
\emph{minimal} typed extension of LangGraph's channel model that is fail-closed, and any
alternative safe reducer on $D$ strictly increases the escalation rate.
\end{theorem}

\begin{proof}
$D$ is a chain; apply Theorem~\ref{thm:general_char} and Corollaries~\ref{cor:chain}
and~\ref{cor:precision}. \qed
\end{proof}

\subsection{Effectiveness: Structural Atomic Gate}
\label{sec:gate}

\begin{proposition}[\texttt{interrupt\_before} addresses workflow coordination, not protocol validity]
\label{prop:interrupt}
LangGraph's \texttt{interrupt\_before} was designed for human-in-the-loop workflow coordination, not fraud-semantic gating. An agent using \texttt{interrupt\_before=["execute\_payment"]} presents graph state to the reviewer---not wallet state, idempotency check results, or action-sequence history. The checkpoint provides no protocol-validity signal. This is a natural extension point: ACP's gate can be composed with \texttt{interrupt\_before} to add fraud-semantic verification at the same workflow step without replacing the coordination mechanism.
\end{proposition}

\textbf{Effectiveness failure in LangGraph.} The following counterexample demonstrates an Effectiveness failure: \texttt{reason\_node} overwrites \texttt{block\_node}'s decision using last-writer-wins semantics, allowing \texttt{pay\_node} to execute without the gate's approval. Under CTS with a \texttt{block\_wins} reducer, this fails at the type level: \texttt{block\_wins("block", "allow") = "block"} and the state cannot decrease.

\begin{lstlisting}
from langgraph.graph import StateGraph
from typing import TypedDict

class State(TypedDict):
    decision: str  # last-writer-wins: no block_wins reducer

def block_node(s): return {"decision": "block"}   # fraud gate fires
def reason_node(s): return {"decision": "allow"}  # overwrites silently
def pay_node(s):
    assert s["decision"] == "allow"               # always passes
    return s  # payment executes despite prior block

g = StateGraph(State)
g.add_node("block", block_node)
g.add_node("reason", reason_node)
g.add_node("pay", pay_node)
g.add_edge("block", "reason"); g.add_edge("reason", "pay")
# result: pay_node executes even though block_node returned "block"
\end{lstlisting}

Under CTS with a \texttt{block\_wins} reducer on the \texttt{decision} channel, the same pipeline would raise a validation error on the overwriting write from \texttt{reason\_node}, because \texttt{block\_wins("block", "allow") = "block"} and the state cannot decrease.

\begin{definition}[Structural Atomic Gate]
A node $G$ in the execution graph is a \emph{structural atomic gate} iff:
\begin{enumerate}
\item[(a)] $G$ executes synchronously before any consequential action reaches the payment rail
\item[(b)] $G$ produces a signed Commerce Action Receipt (CAR) with a \texttt{final\_decision} field
\item[(c)] The graph has no path from $G$ to the payment rail bypassing \texttt{final\_decision}
\item[(d)] $G$'s output is persisted to the tamper-evident CARLog before the payment rail executes
\end{enumerate}
\end{definition}

Properties (a)--(d) formalize the 2PC coordinator role: $G$ is the coordinator, the payment rail is the participant. No participant commits without coordinator approval; the decision is durably logged before participant execution.

\begin{definition}[Commerce Action Receipt (CAR)]
A CAR is a tuple $(\textit{id}, \textit{action}, \textit{verifications}, \textit{decision}, \textit{level}, \textit{anchoring}, \textit{sig}, \textit{ts})$ where \textit{verifications} lists $(\textit{verifier\_id}, \textit{score}, \textit{confidence}, \textit{uncertainty})$; $\textit{decision} \in \{\textit{allow}, \textit{flag}, \textit{block}\}$; $\textit{level} \in \{L_0, L_1, L_2, L_3\}$; \textit{anchoring} $= (\textit{log\_hash}, \textit{prev\_hash}, \textit{merkle\_root}, \textit{chain\_tx})$; \textit{sig} $= \text{HMAC-SHA256}$ over decision-relevant fields.
\end{definition}

\textbf{Implementation conformance.} The ACP ceremony node is designed to satisfy Definition~3.4: (a)~ceremony routing is synchronous ($<$10ms); (b)~CAR is signed before payment routing; (c)~routing is conditional on \texttt{final\_decision}, so no path exists to the payment rail bypassing the gate; (d)~CAR is appended to the hash-chained CARLog before the payment node executes. Together with Theorems~\ref{thm:type_safety}--\ref{thm:char}, this gives the Effectiveness property: a block from any verifier tier cannot be overwritten by downstream nodes.

\subsection{Atomicity: Risk-Adjusted Commitment Protocols}
\label{sec:tiers}

\begin{table}[h]
\caption{Ceremony tiers: risk-adjusted commitment protocols. $L_0$--$L_2$ implemented; $L_3$ spec only (see \S\ref{sec:discussion_limits}).}
\label{tab:tiers}
\centering
\footnotesize
\begin{tabular}{lp{3.5cm}lll}
\toprule
Level & Condition & Protocol & Anchoring & Status \\
\midrule
$L_0$ & \texttt{FIND}, \texttt{QUOTE} & None & None & Impl. \\
$L_1$ & \texttt{COMMIT}, amt$<$\$50, unc$<$0.25 & Saga & Local hash-chain & Impl. \\
$L_2$ & \texttt{COMMIT}, \$50--\$5k, or unc$\geq$0.25 & 2PC-lite & Merkle batch & Impl. \\
$L_3$ & amt$>$\$5k, unc$\geq$0.45, or disagreement & 2PC+consensus & On-chain & Spec only \\
\bottomrule
\end{tabular}
\end{table}

The ceremony tier assignment is risk-adjusted: low-stakes actions (FIND, QUOTE) require no commitment protocol because they produce no irreversible state change. COMMIT actions below \$50 with low uncertainty follow the Saga pattern---each action carries a compensating transaction (VOID) that can be executed on failure. Higher-value or higher-uncertainty actions escalate to 2PC-lite, where the CAR coordinator must receive prepare acknowledgment from the payment participant before issuing commit. A RESERVE-then-failure at any tier triggers automatic compensating VOID, satisfying the Atomicity property by construction.

\subsection{Traceability: Hash-Chained CARLog}
\label{sec:carlog}

Adapting Certificate Transparency~\cite{ct}:
\begin{align*}
\textit{entry}_n = \{\ &\textit{seq}: n,\ \textit{car\_hash}: \text{SHA256}(\textit{CAR}),\\
&\textit{prev\_hash}: \text{SHA256}(\textit{entry}_{n-1}),\ \textit{entry\_sig}: \text{HMAC}(\cdot)\ \}
\end{align*}

\begin{proposition}[Tamper-evidence]
Any modification of $\textit{entry}_k$ changes $\textit{entry}_k.\textit{entry\_hash}$, which changes $\textit{entry}_{k+1}.\textit{prev\_hash}$, breaking all subsequent entry signatures. A chain scan detects modification in $O(N)$ time.
\end{proposition}

Merkle batch anchoring: every $B$ CARs are batched into a Merkle tree; the root is written to a public blockchain. An inclusion proof for $\text{CAR}_k$ is $O(\log B)$ sibling hashes---verifiable by any third party without trusting the operator, and bound to a specific block timestamp.

\subsection{Verification Depth: Confidence Cascade}
\label{sec:cascade}

The cascade implements selective prediction~\cite{selective}: a verifier uncertain beyond threshold $\theta$ should abstain and defer to a higher tier.

\begin{lstlisting}
Tier 1 (LocalRulesVerifier):  confidence = 1.0, uncertainty = 0.0
   -> rule fires: block, no delegation
   -> rule passes: delegate to Tier 2

Tier 2 (BehavioralML):        confidence = 1 - uncertainty_heuristic
   -> uncertainty: 0.30 if |score - threshold| < 0.12, else 0.08
   -> rules-certainty veto: downgrade block->flag if LocalRules score=0 AND unc>0.20

Tier 3 (LLM Verifier):        invoked for <3% of actions (high uncertainty)
\end{lstlisting}

Note: current uncertainty estimates are \emph{heuristic} (boundary proximity, OOD boost)---not true conformal prediction intervals~\cite{conformal}. Production deployment requires a held-out calibration set from the deployment distribution before activating conformal coverage guarantees.

\begin{proposition}[Cascade monotonicity]
Under the CTS \texttt{block\_wins} reducer, adding Tier~2 or Tier~3 cannot reduce a block from a prior tier, and can only add flags or blocks. The cascade is monotone in the safety direction.
\end{proposition}

\section{Experiments}
\label{sec:experiments}

We measure all five transaction properties across 900 synthetic commerce sessions (seeds 42--46, 100 sessions per attack type plus 100 clean sessions per seed). The companion benchmark paper~\cite{acpbench} describes the full data generation protocol.

\subsection{Sanity Check: Structural Property Coverage}

Before measuring the research contribution, we verify that CAO satisfies the five EAVT properties in the baseline sense: the $\mathcal{M}$-check accepts valid sessions and rejects grammar violations (E); all COMMIT/RESERVE actions pass through the gate (Ef); CARLog records every action in a hash-chained audit trail (T); Atomicity and Verification hold by architectural construction. These are expected from the design---not research results. Naive execution satisfies none; schema-only achieves syntax-level Ef only. This comparison is included as a correctness baseline, not a claim.

Specifically: CAO catches 300/300 definitional M-violation events (B1, B7, B8-clear) and gates 100\% of 2010 COMMIT/RESERVE actions. The $\mathcal{M}$-check is a DFA simulation---it is definitionally correct on attacks that violate the grammar. The M-check cannot detect behavioral attacks (B2--B6) crafted to satisfy all grammar transitions and field constraints; those are the research problem handled by the confidence cascade.

\subsection{Formal Property Verification}
\label{sec:formal_props}

\textbf{Safety monotonicity} (Theorem~\ref{thm:monotone}): in 1,000 random session orderings (permuting verifier execution order), session-level decisions are identical in 100\% of cases. The \texttt{block\_wins} reducer is order-independent by construction.

\textbf{Tamper-evidence}: random byte flips at 10 positions in a 172-entry CARLog. \texttt{log.verify()} detects 10/10 corruptions in $O(N)$ time.

\textbf{Inclusion proof validity}: Merkle inclusion proofs verified for 172/172 entries.

\subsection{Fraud Detection as Measured Consequence}

A system that enforces Execution Correctness and Verification is also structurally resistant to protocol-layer attacks. We verify this on ACP-Bench~\cite{acpbench} and report the results as a secondary finding. The fraud detection improvement is not a separate design goal---it is a measured consequence of satisfying the EAVT properties.

\begin{table}[h]
\caption{ACP-Bench results: ML only and ACP rows are mean $\pm$ std across seeds 42--46 over the core 8-attack benchmark (520 sessions); single-seed reference (seed=42) in parentheses. The LLM-judge row is seed=42, all 10 attacks (measured). \emph{Behavioral recall} (B2--B6, B7s, B8-marginal/B8s) is the research metric; structural attack recall (B1/B7/B8-clear) is expected from the DFA check and is not the contribution. B7s and B8s are behavioral variants that rules catch at 0\% by construction, confirming the benchmark requires behavioral modeling beyond rule checking. \textit{Total FPR} (block+flag on clean) is identical for ACP and MLOnly at every seed.}
\label{tab:results}
\centering
\footnotesize
\begin{tabular}{lccccc}
\toprule
Approach & F1 & Behavioral Recall & Hard FPR & EscRate & \$/1k \\
\midrule
No verification & 0.000 & 0.000 & 0.000 & 0.000 & $<$\$0.01 \\
Schema only (LangGraph) & 0.000 & 0.000 & 0.000 & 0.000 & $<$\$0.01 \\
Rules only & $0.573{\pm}0.008$ & $0.000$ & $\mathbf{0.000}$ & $\mathbf{0.000}$ & $<$\$0.01 \\
 & \small(0.584) & \small(0.000) & & & \\[2pt]
LLM judge, Sonnet (L0 ceil.) & n/a & 0.610 & 0.000 & 0.000 & \$3.00 \\[2pt]
ML only & $0.712{\pm}0.012$ & $0.740{\pm}0.021$ & $0.090{\pm}0.028$ & $0.504{\pm}0.034$ & $<$\$0.01 \\
 & \small(0.712) & \small(0.747) & \small(0.110) & \small(0.450) & \\[2pt]
\textbf{ACP (CAO)} & $\mathbf{0.792{\pm}0.006}$ & $\mathbf{0.880{\pm}0.009}$ & $\mathbf{0.000{\pm}0.000}$ & $0.594{\pm}0.028$ & $<$\$0.01 \\
 & \small(0.792) & \small(0.884) & \small(\textbf{0.000}) & \small(0.560) & \\
\bottomrule
\end{tabular}
\end{table}

\textbf{The research result is not the rules column.} Rules-only achieves 0\% behavioral recall (B2--B6, B7s, B8-marginal/B8s require behavioral signal absent from the rule checker) and 0\% hard FPR---these are expected outcomes from a rule checker, not findings. B7s and B8s specifically confirm this: rules achieve 0\% on both behavioral variants by construction (the MCC is allowlisted; the individual amounts are within ceilings), yet MLOnly achieves 82.5\% on B7s and 100\% on B8s. The L0 ceiling---Claude Sonnet~4.6 full-session judge (measured, not approximated)---achieves 61\% recall at 0\% hard FPR: strong on text-visible semantics (B6/B2/B8s at 100\%, B7s at 95\%) but limited by absent L1 state (B4/B5/B8 at 10--25\%). \textbf{The finding} is the cascade composition: ML-only achieves 74\% behavioral recall at 9\% hard FPR; ACP achieves 88\% behavioral recall at 0\% hard FPR, surpassing the strongest L0 ceiling by 27 recall percentage points at 300$\times$ lower cost. The $\Delta$ (recall +14 percentage points, hard FPR --9 percentage points) comes from two mechanisms: (1) rules provide coverage for the structural cases ML misses due to behavioral feature blindness, and (2) the rules-certainty veto converts ML's uncertain blocks on high-value legitimate sessions to auditable flags rather than auto-rejections. Total FPR is identical for ACP and MLOnly at every seed---the operational benefit is elimination of auto-blocks, not reduction of review load.

\section{Discussion}

\subsection{Theoretical Relationship to Prior Work}

CAO's five properties map directly to classical distributed systems guarantees: Execution correctness corresponds to ISO~8583 grammar enforcement, which terminal hardware provided implicitly for four decades---CAO makes this explicit and applies it to LLM agents. Atomicity corresponds to the Saga/2PC guarantee that partial execution produces compensating rollback. Effectiveness corresponds to the 2PC coordinator role, which ensures no participant commits without coordinator approval. Verification corresponds to selective prediction, where uncertain verifiers abstain and defer to higher tiers. Traceability corresponds to Certificate Transparency, providing an append-only, hash-chained, Merkle-anchored audit log verifiable by third parties without trusting the operator.

\textbf{CAO vs.\ Saga}: Saga says ``how to undo.'' CAO says ``whether to proceed''---the ceremony gate must approve before any action is committed, even one with a compensating transaction.

\textbf{CAO vs.\ 2PC}: 2PC provides consensus among participants. CAO adds a pre-commit verification gate: the agent verifier acts as coordinator, approving before the payment rail (participant) executes its own 2PC. CAO does not replace 2PC in the rail---it layers above it.

\textbf{CTS vs.\ LangGraph state}: LangGraph reducers are unconstrained. CTS constrains reducers to monotone semilattice joins. Theorem~\ref{thm:lww_unsound} proves LWW is \emph{necessarily} unsafe; Theorem~\ref{thm:char} proves that monotone joins with \texttt{block} as top are the minimal safe typing discipline---CTS is not one option among many but the unique minimal fix. This is a strict strengthening of LangGraph's model, not merely an alternative design.

\textbf{$\mathcal{M}$ vs.\ prior state machine work}: the commerce state automaton is not a new theoretical construct---it formalizes what ISO~8583 has enforced implicitly for decades. The novelty is applying it as a security invariant for agentic systems where no terminal hardware enforces the grammar.

\subsection{Limitations}
\label{sec:discussion_limits}

\textbf{Heuristic uncertainty, not conformal prediction.} The BehavioralML uncertainty estimates are boundary-proximity heuristics: $\hat{u} = 0.30$ if $|\text{score} - \text{threshold}| < 0.12$, else $0.08$, with additive boosts for cold-start and OOD patterns. These are \emph{not} conformal-prediction intervals~\cite{conformal}: (a)~they carry no marginal coverage guarantee on a held-out calibration set from the deployment distribution; (b)~they do not account for covariate shift between training (seed=7, $N=500$) and deployment; (c)~the veto threshold (uncertainty~$> 0.20$) and escalation threshold (uncertainty~$> 0.35$) are empirically tuned, not derived from a coverage level. The formal framework in \S\ref{sec:cts} assumes uncertainty quantification that supports the selective-prediction guarantee~\cite{selective}---the current implementation does not fully realize this. Closing this gap requires: (i)~a calibration set from the target deployment distribution; (ii)~replacing $\hat{u}$ with a proper non-conformity score; (iii)~selecting thresholds at a pre-specified $1-\alpha$ coverage level before evaluation.

\textbf{Escalation rate (56\%) is mechanistically explained, not a calibration artifact.} With 0\% hard FPR, all clean-session false positives are soft flags (EscRate=56\%). These are sessions where ML scores above threshold due to high-value travel amounts (\$1,800--\$2,800), but LocalRules score$=0$ (below the \$3,000 ceiling), so the rules-certainty veto (Proposition~\ref{prop:veto}) correctly downgrades block$\to$flag. Reducing EscRate without sacrificing B8 recall requires per-persona amount thresholds---future work.

\textbf{Incomplete attack classes.} B9 (revert-grant) requires mocking post-settlement rail state and is deferred. B4 (session hijack) at 68\% requires an external agent-ID registry. B5 (idempotency replay) at 70\% is via ML behavioral anomaly; the architecturally correct mechanism is a persistent settled-key registry in the MCP server. B7s (MCC spoof) at 82.5\% requires merchant-domain knowledge (a merchant taxonomy mapping domain to risk category) beyond the current wallet allowlist check.

\textbf{L3 consensus not yet implemented.} The CAR and CARLog infrastructure is implemented and formally specified, but the Byzantine-fault-tolerant multi-verifier consensus protocol for $L_3$ ceremony actions is not yet implemented. The formal specification is in \S\ref{sec:gate}; production realization requires a quorum protocol adapted for the verifier-coordinator role.

\section{Conclusion}

Existing agentic orchestration frameworks provide none of the five transaction properties---Execution correctness, Atomicity, Effectiveness, Verification, and Traceability---that sound commerce systems require. We have shown that CAO provides all five, and proved three characterization results that explain \emph{why} existing frameworks cannot. CTS Type Safety (Theorem~\ref{thm:type_safety}) proves that a correctly-typed CTS graph cannot bypass a prior block decision. LWW Unsoundness (Theorem~\ref{thm:lww_unsound}) proves that any LangGraph-compatible graph using last-writer-wins \emph{must} admit such a bypass---it is a structural impossibility, not a configuration issue. The Characterization Theorem (Theorem~\ref{thm:char}) proves that monotone joins with \texttt{block} as top are the \emph{unique minimal} safe reducer semantics, making CTS the minimal typed extension of LangGraph that achieves fail-closed safety.

The grammar enforcement ($\mathcal{M}$-check) and audit logging (CARLog) are correct by construction---expected, not novel research results. The research contribution is the characterization trilogy and the cascade composition that operationalizes it: the L0 ceiling (Claude Sonnet full-session LLM judge) achieves 61\% recall at \$3/1k; ML-only achieves 74\% at 9\% hard FPR; ACP achieves 88\% at 0\% hard FPR at $<$\$0.01/1k---surpassing the L0 ceiling by 27 percentage points at 300$\times$ lower cost. B4/B5 at 68--70\% remain open; closing them requires cross-session infrastructure beyond single-session typed state.

The transition from semantically neutral graph execution to commerce-typed execution with invariant-enforcing reducers is a qualitative shift in what can be formally guaranteed about agent behavior. As agents are authorized to take increasingly consequential financial actions, the theoretical foundations developed here become not optional but necessary.

\section*{Broader Impact}

This framework enables static and runtime compliance tooling for agentic payment systems, potentially reducing financial fraud while improving precision (fewer false-positive blocks). Publishing the automaton grammar formalizes what ISO~8583 has enforced for four decades, so it does not constitute novel attack information; the characterization theorem (Theorem~\ref{thm:char}) additionally proves that CTS is the minimal typed extension needed to close the bypass---adversaries cannot exploit a weaker typing discipline to circumvent the gate. The released artifact is the verifier, not exploit code.

\begin{ack}
This work was conducted without external funding. Competing interests: none.
\end{ack}

\section*{Reproducibility Statement}

Formal results (Characterization Theorem~\ref{thm:general_char} with full proof in Appendix~\ref{app:general_char}, Delegation Closure Theorem~\ref{thm:delegation}, Corollaries~\ref{cor:chain}--\ref{cor:precision}, CTS Type Safety Theorem~\ref{thm:type_safety}, LWW Unsoundness Theorem~\ref{thm:lww_unsound}, commerce instantiation Theorem~\ref{thm:char}, Propositions~\ref{prop:decidability}--\ref{prop:veto}) are self-contained in the text and appendix. The LangGraph counterexample (Proposition~\ref{prop:interrupt}) is 15 lines of Python reproducible with \texttt{pip install langgraph}. All empirical results are reproducible via the companion benchmark~\cite{acpbench}: \texttt{python -m benchmark.run\_experiment --seed \{42..46\} --no-sota}. Implementation details: $L_0$--$L_2$ ceremony tiers, hash-chained CARLog, and Merkle batch anchoring are implemented; $L_3$ consensus is a formal specification only.

\begin{thebibliography}{99}
\small

\bibitem{saga} Garcia-Molina, H., \& Salem, K. (1987). Sagas. \emph{ACM SIGMOD Record}, 16(3), 249--259.

\bibitem{2pc} Gray, J.\ N. (1978). Notes on data base operating systems. \emph{Operating Systems: An Advanced Course}, 393--481.

\bibitem{iso8583} ISO/IEC 8583:2021. \emph{Financial transaction card originated messages: Interchange message specifications}. ISO.

\bibitem{ct} Laurie, B., Langley, A., \& Kasper, E. (2013). Certificate Transparency. \emph{RFC~6962}. IETF.

\bibitem{react} Yao, S., et al.\ (2022). ReAct: Synergizing reasoning and acting in language models. \emph{ICLR 2023}.

\bibitem{langgraph} LangChain Inc.\ LangGraph. \url{https://langchain-ai.github.io/langgraph/}, 2024.

\bibitem{pregel} Malewicz, G., et al.\ (2010). Pregel: A system for large-scale graph processing. \emph{ACM SIGMOD}, 135--146.

\bibitem{autogen} Wu, Q., et al.\ (2023). AutoGen: Enabling next-gen LLM applications via multi-agent conversation. \emph{arXiv:2308.08155}.

\bibitem{openai_agents} OpenAI. OpenAI Agents SDK. \url{https://platform.openai.com/docs/guides/agents}, 2025.

\bibitem{x402} Li, et al.\ ``Security Analysis of the x402 AI Payment Protocol.'' \emph{arXiv:2605.11781}, 2026.

\bibitem{crdts} Shapiro, M., et al.\ (2011). Conflict-free replicated data types. \emph{SSS 2011}, 386--400.

\bibitem{typestate} Strom, R.\ E., \& Yemini, S. (1986). Typestate: A programming language concept for enhancing software reliability. \emph{IEEE Trans.\ Software Eng.}, 12(1), 157--171.

\bibitem{sessiontypes} Honda, K., Vasconcelos, V.\ T., \& Kubo, M. (1998). Language primitives and type discipline for structured communication-based programming. \emph{ESOP 1998}, 122--138.

\bibitem{selective} Geifman, Y., \& El-Yaniv, R. (2017). Selective prediction in deep neural networks. \emph{NeurIPS 2017}.

\bibitem{conformal} Angelopoulos, A.\ N., \& Bates, S. (2022). A gentle introduction to conformal prediction. \emph{arXiv:2107.07511}.

\bibitem{stripe} Stripe Inc.\ ``PaymentIntent object.'' \url{https://stripe.com/docs/api/payment_intents}, 2024.

\bibitem{acpbench} [Authors]. ACP-Bench: Measuring Safety at the Payment-Protocol Layer of Agentic Commerce. \emph{AIWILD @ NeurIPS 2026}. (companion paper)


\end{thebibliography}

\appendix

\section{Formal Proofs}

\subsection{Proof of Theorem~\ref{thm:general_char} (Characterization of safe reducers)}
\label{app:general_char}

Throughout, $r$ is a channel reducer on a bounded safety lattice $L$, so $r(\bot,x)=x$, and
$D_i$ denotes the fold of a trace $b_1,\dots,b_n$.

\textbf{($\Leftarrow$).} Assume $r$ is associative, commutative, idempotent, and $x\sqcup y\le r(x,y)$.
(S3): associativity and commutativity make the fold of any trace equal to the $r$-product of its
multiset, which is permutation-invariant. (S4): idempotence makes the $r$-product of a multiset
equal to that of its underlying set, so duplicates do not change $D_n$. (S2): $D_{i-1}\le
D_{i-1}\sqcup b_i\le r(D_{i-1},b_i)=D_i$ and likewise $b_i\le D_i$. (S1): if $b_i=\top$ then by
(S2) $\top\le D_i$, so $D_i=\top$; by (S2) again $D_j\ge D_i=\top$ for all $j\ge i$.

\textbf{($\Rightarrow$).} Assume (S1)--(S4). Every $x\in L$ is a reachable state (trace $(x)$ gives
$D_1=r(\bot,x)=x$), so trace conditions yield identities on all of $L$.
\emph{Commutativity}: traces $(x,y)$ and $(y,x)$ give $D_2=r(x,y)$ and $r(y,x)$; (S3) equates them.
\emph{Associativity}: traces $(x,y,z)$ and $(y,z,x)$ give $r(r(x,y),z)$ and $r(r(y,z),x)=r(x,r(y,z))$
(by commutativity); (S3) equates them.
\emph{Idempotence}: traces $(x)$ and $(x,x)$ give $x$ and $r(x,x)$; (S4) equates them.
\emph{Domination}: (S2) on trace $(x,y)$ gives $x=D_1\le D_2=r(x,y)$ and $y=b_2\le D_2=r(x,y)$; so
$r(x,y)$ is an upper bound of $\{x,y\}$ in $(L,\le)$ and therefore $x\sqcup y\le r(x,y)$.

\textbf{Minimality.} $\sqcup$ is a semilattice operation dominating itself, hence safe by
($\Leftarrow$). For any safe $r$, ($\Rightarrow$) gives $\sqcup\le r$ pointwise. If $r\ne\sqcup$
then $r(x,y)>x\sqcup y$ for some pair, which is a strict escalation over $\sqcup$ on that pair
with no gain in (S1)--(S4). \qed

\subsection{Proof of Proposition (CTS channels are CRDTs)}

The \texttt{block\_wins} reducer is a join-semilattice on the lattice $\texttt{allow} < \texttt{flag} < \texttt{block}$: associative (three-element total order is associative), commutative (lattice join is symmetric), and idempotent ($f(x,x) = x$ for all $x$). The \texttt{operator.add} reducer over \texttt{Decimal} is a join-semilattice on non-negative reals with standard $\leq$. The \texttt{frozenset.union} reducer is a join-semilattice on the powerset lattice with subset ordering. All three are CRDTs by the CRDT join-semilattice characterization~\cite{crdts}.

\subsection{Proof of Proposition~\ref{prop:veto} (Rules-Certainty Veto)}

Let $b_i \in D = \{\texttt{allow}, \texttt{flag}, \texttt{block}\}$ be the \emph{action-level} output of the ACP node at step $i$ (after veto adjustment if applicable). Let $D_i = \texttt{block\_wins}(D_{i-1}, b_i)$ be the session-level decision after step $i$ with $D_0 = \texttt{allow}$.

\textbf{Veto trigger}: fires at step $i$ when (1)~LocalRules score$=0$ (no deterministic violation detected for action $i$); (2)~ML decision$_{\mathrm{raw}} = \texttt{block}$; (3)~uncertainty $> 0.20$. The veto adjusts $b_i \leftarrow \texttt{flag}$ before the session-level reducer runs.

\textbf{Key observation}: condition (1) means LocalRules produced decision$=\texttt{allow}$ for action $i$ specifically. It says nothing about prior actions $j < i$. The session-level state $D_{i-1}$ may already be \texttt{block} from a prior step.

\textbf{Case 1} ($D_{i-1} = \texttt{block}$): $D_i = \texttt{block\_wins}(\texttt{block}, \texttt{flag}) = \texttt{block}$. The veto cannot reduce a prior block. (Follows from Theorem~\ref{thm:type_safety}.)

\textbf{Case 2} ($D_{i-1} < \texttt{block}$): $D_i = \texttt{block\_wins}(D_{i-1}, \texttt{flag}) = \texttt{flag}$ (since \texttt{flag} $>$ \texttt{allow} in the order and $D_{i-1} \in \{\texttt{allow}, \texttt{flag}\}$). The veto correctly converts an uncertain ML block to a reviewable flag without triggering an auto-block.

In both cases the veto does not violate Theorem~\ref{thm:type_safety}: it operates at the action level before the session-level join, and the join remains monotone. \qed

\section{Extensions}

\textbf{A2A delegation and sub-agent scope.} As agents delegate to sub-agents, the CAR provenance model should track the full delegation chain. A sub-agent's \texttt{COMMIT} should be CARed under the delegating agent's wallet with explicit scope constraints.

\textbf{Financial trading and crypto.} The automaton $\mathcal{M}$ and CTS framework are asset-class agnostic. Equity trading agents have an analogous lifecycle (\texttt{QUOTE\_PRICE}$\to$\texttt{PLACE\_ORDER}$\to$\texttt{FILL}$\to$\texttt{CONFIRM}$\to$\texttt{SETTLE}) with a different field-constraint vocabulary (instrument allowlist, position size, leverage ratio). Crypto agents require \texttt{Reversibility=NONE} at \texttt{COMMIT} time and a chain-ID field for cross-chain replay prevention.

\textbf{Observer mode.} An observer-mode extension produces risk scores and chargeback-evidence packages without inline interception, for deployments where intermediation is not architecturally feasible.

\section*{NeurIPS Paper Checklist}

\begin{enumerate}

\item {\bf Claims}
    \item[] Question: Do the main claims made in the abstract and introduction accurately reflect the paper's contributions and scope?
    \item[] Answer: \answerYes{}
    \item[] Justification: The abstract distinguishes two types of claims: (1) structural coverage metrics (100\% M-check catch rate, FPR~=~0\% by DFA construction; 100\% CARLog traceability, architectural guarantee) and (2) probabilistic fraud detection metrics (88.0\%$\pm$0.9\% recall, 0.0\%$\pm$0.0\% hard FPR, 59.4\% escalation rate --- multi-seed over 8-attack benchmark, matched exactly in Table~\ref{tab:results}). B7s/B8s per-attack results and the LLM-judge baseline (61\% recall) are seed=42 supplementary results, labeled as such in \S6. The distinction between structural coverage claims and probabilistic detection metrics is elaborated in \S6. The Commerce State Automaton is formally defined (Definition~3.0) with proofs provided for all three theoretical results. No L3 implementation is claimed.

\item {\bf Limitations}
    \item[] Question: Does the paper discuss the limitations of the work performed by the authors?
    \item[] Answer: \answerYes{}
    \item[] Justification: \S\ref{sec:discussion_limits} covers: (1) heuristic uncertainty estimates are boundary-proximity heuristics, not conformal-prediction intervals---the exact conditions required for marginal coverage are stated; (2) 56\% escalation rate (soft-flag FPR; hard FPR is 0.0\%) is a per-persona spend-ceiling overlap limitation, not a calibration artifact; (3) B7s (MCC spoof) at 82.5\% requires merchant-domain knowledge beyond the current allowlist; (4) B4/B5 at 68--70\% require cross-session registries; (5) L3 ceremony tier (ZK proof, selective-disclosure credential) is unimplemented; (6) the commerce automaton is validated against ISO~8583, x402, and Stripe PaymentIntent but not against other payment networks.

\item {\bf Theory assumptions and proofs}
    \item[] Question: For each theoretical result, does the paper provide the full set of assumptions and a complete proof?
    \item[] Answer: \answerYes{}
    \item[] Justification: The main results are stated with explicit assumptions and complete proofs. Theorem~\ref{thm:general_char} (characterization of safe reducers on any finite bounded lattice: safe iff semilattice operation dominating the join; join is the unique minimal safe reducer) has a proof sketch in \S\ref{sec:general_typing} and a full proof in Appendix~\ref{app:general_char}; its assumptions (finite bounded lattice, reducer with $\bot$ as unit, trace semantics S1--S4) are stated in Definitions preceding it. Theorem~\ref{thm:delegation} (delegation closure: composite is fail-closed iff the lattice embedding is top-preserving) is proved in full in \S\ref{sec:general_typing}. Theorems~\ref{thm:type_safety}, \ref{thm:lww_unsound}, \ref{thm:char} (commerce instantiation) are proved in \S\ref{sec:cts}. Proposition~\ref{prop:decidability} (decidability, $O(n)$) and the tamper-evidence and CRDT propositions are proved in the text and appendix. Where a result is a direct specialization (e.g., Theorem~\ref{thm:char} from Theorem~\ref{thm:general_char}), the paper says so explicitly.

\item {\bf Experimental result reproducibility}
    \item[] Question: Does the paper fully disclose all information needed to reproduce the main experimental results?
    \item[] Answer: \answerYes{}
    \item[] Justification: The Reproducibility Statement gives the exact command (\texttt{python -m benchmark.run\_experiment --seed 42 --no-sota}), compute environment (CPU only, $<$30 seconds/seed), and JSON output field names. Markov parameters, amount distributions, perturbation $\sigma$, training seed, threshold formula, and veto conditions are fully specified in the companion AIWILD paper (cited), which contains the benchmark details. All automaton transitions are fully enumerated in Table~1.

\item {\bf Open access to data and code}
    \item[] Question: Does the paper provide open access to data and code?
    \item[] Answer: \answerYes{}
    \item[] Justification: ACP reference implementation and benchmark harness are released as open-source (anonymized URL in supplemental material; full URL in camera-ready). The benchmark is fully deterministic given \texttt{--seed}. The formal automaton definition (Definition~3.0) and CTS channel specification (\S4) are fully self-contained and implementation-independent.

\item {\bf Experimental setting/details}
    \item[] Question: Does the paper specify all training and test details necessary to understand the results?
    \item[] Answer: \answerYes{}
    \item[] Justification: Training data: seed=7, $N=500$ sessions. Test data: seeds 42--46, $N=520$ sessions each. ML threshold: $\mu + 1.1\sigma$ of clean training scores. Veto: score$=0$ and uncertainty$>0.20$. Uncertainty: boundary-proximity heuristic (0.30 if $|$score$-$threshold$|<0.12$, else 0.08). All specified in companion AIWILD paper (\S3 and Appendix~A).

\item {\bf Experiment statistical significance}
    \item[] Question: Does the paper report error bars or appropriate information about statistical significance?
    \item[] Answer: \answerYes{}
    \item[] Justification: All multi-seed results are reported as mean $\pm$ standard deviation across 5 independent seeds (42--46). Standard deviations are population standard deviations, not standard errors---stated in the table caption. The five seeds are fully independent (separate data generation). Rates are bounded well away from 0 and 1 in the reported range, so symmetric intervals are appropriate.

\item {\bf Experiments compute resources}
    \item[] Question: Does the paper provide sufficient information on compute resources needed to reproduce the experiments?
    \item[] Answer: \answerYes{}
    \item[] Justification: CPU only; no GPU required. $<$30 seconds per seed; $<$3 minutes for all five seeds. The Tier~3 LLM Verifier invoked on $<$3\% of actions in the full pipeline is bypassed by \texttt{--no-sota} in all reported experiments. Stated in the Reproducibility Statement.

\item {\bf Code of ethics}
    \item[] Question: Does the research conform with the NeurIPS Code of Ethics?
    \item[] Answer: \answerYes{}
    \item[] Justification: This work provides formal safety guarantees and defensive tooling for financial AI systems. All benchmark data is synthetic. No production payment credentials or working exploit code are released. The attack taxonomy is presented at the conceptual level required to motivate the defensive framework.

\item {\bf Broader impacts}
    \item[] Question: Does the paper discuss both potential positive and negative societal impacts?
    \item[] Answer: \answerYes{}
    \item[] Justification: See the Broader Impact section following the Conclusion. Positive: formal safety guarantees at the protocol layer (L1) reduce the window for financial fraud in agentic systems; the decidability result enables static compliance tooling. Negative: publication of the automaton model could inform adversarial agents that probe state transition boundaries; mitigated by the verifier being released alongside the taxonomy.

\item {\bf Safeguards}
    \item[] Question: Does the paper describe safeguards for responsible release of data or models with high risk for misuse?
    \item[] Answer: \answerYes{}
    \item[] Justification: The formal framework defines safety properties; the release includes the verifier (defensive tool) and the synthetic generator, not production payment attack implementations. No live API credentials, no production merchant exploit code, and no B5 settled-key replay implementations are released.

\item {\bf Licenses for existing assets}
    \item[] Question: Are creators of assets used in the paper properly credited and license terms mentioned?
    \item[] Answer: \answerYes{}
    \item[] Justification: LangGraph~\cite{langgraph} (Apache 2.0); ISO~8583 is a published standard (cited); x402~\cite{x402} (open protocol, arXiv); Stripe~\cite{stripe} (public API documentation). All appear in the References section.

\item {\bf New assets}
    \item[] Question: Are new assets introduced in the paper well documented?
    \item[] Answer: \answerYes{}
    \item[] Justification: The Commerce State Automaton (Definition~3.0), CTS channel type (Definition~4.1), and CARLog (Definition~4.2) are formally defined with full enumeration of states, transitions, and operators. The ACP reference implementation is documented in \S4 and released open-source under MIT license.

\item {\bf Crowdsourcing and research with human subjects}
    \item[] Question: Does the paper include full details of crowdsourcing or human subjects research?
    \item[] Answer: \answerNA{}
    \item[] Justification: All experimental data is synthetically generated. No human subjects, no crowdsourcing, no real user data.

\item {\bf Institutional review board (IRB) approvals}
    \item[] Question: Does the paper describe potential risks to study participants and IRB approvals?
    \item[] Answer: \answerNA{}
    \item[] Justification: No human subjects research. Not applicable.

\item {\bf Declaration of LLM usage}
    \item[] Question: Does the paper describe the usage of LLMs if it is an important, original, or non-standard component of the core methods?
    \item[] Answer: \answerYes{}
    \item[] Justification: Two LLM uses are disclosed. (1) The Tier~3 LLM Verifier (\S4) uses \texttt{claude-haiku-4-5-20251001} for high-uncertainty actions ($<3\%$ of actions), described in \S4 (``Ceremony Tiers and Escalation Protocol''). (2) The full-session LLM-as-judge baseline (\S6, Table~\ref{tab:results}) uses \texttt{claude-sonnet-4-6} via AWS Bedrock to establish the L0 detection ceiling (61\% recall, seed=42)---an evaluation component, not a method component. LLMs are \emph{not} used in the formal automaton, the CTS channel type, or the CARLog---those are provably correct by construction. LLMs were also used for writing assistance in preparing this manuscript, which does not require declaration per the NeurIPS handbook.

\end{enumerate}


\end{document}
